\documentclass[10pt,a4paper]{amsart}
\usepackage[margin=19mm]{geometry}
\usepackage{amsmath,amssymb,mathtools}
\usepackage[scr=rsfs,cal=euler]{mathalfa}
\usepackage{xfrac}
\usepackage[unicode,hidelinks,bookmarks=false]{hyperref}
\newcommand{\ie}{i.e.\ }
\newcommand{\cf}{cf.\ }
\newcommand{\resp}{respectively\ }
\newcommand{\Subsection}{\subsection}
\providecommand{\nobreakdash}{\nobreak\hbox{-}}
\setlength{\emergencystretch}{2em}
\multlinegap0pt
\usepackage[shortlabels]{enumitem}
\usepackage{amssymb}
\usepackage{mathtools}
\usepackage{bm}
\usepackage{xcolor}
\newtheorem{theorem}{Theorem}
\title{The Non-Linearity Perturbation Threshold: Width Scaling and Landscape Bifurcations in Deep Learning}
\author{Michael Alexander}
\date{Source: arXiv:2604.01334v1 (CC BY 4.0)}
\begin{document}
\maketitle

\noindent\emph{Source and licence.} The Non-Linearity Perturbation Threshold: Width Scaling and Landscape Bifurcations in Deep Learning. Version arXiv:2604.01334v1 (CC BY 4.0), distributed by arXiv under the Creative Commons Attribution 4.0 International licence, http://creativecommons.org/licenses/by/4.0/, as recorded in the arXiv licence metadata for this exact version and shown on the abstract page https://arxiv.org/abs/2604.01334v1. Full licence notice: \texttt{licenses/arxiv-2604.01334-CC-BY-4.0.txt}.\par\smallskip

\noindent\textit{Editorial scope.} Counted unit: the article's theorem thm:near-pitchfork-structure (source0.tex lines 555--575). The article's complete original proof of it is delivered with it (source0.tex lines 577--599, one \texttt{proof} environment of 23 lines). No mathematics is written, repaired or re-derived: the statement and the proof are the article's own lines, in the article's own notation, and no retained line is substituted or re-flowed.  Retained uncounted as the source-local context the proof needs: lines 537--540.  Nothing was omitted from any retained span, so the package carries no omission record.  No retained line contains a citation, so the wrapper carries no bibliography and no bibliography entry.  No retained line contains a figure environment, an image-inclusion command or drawing code, so no image is delivered and the package declares no asset dependency.  Editorial formatting only, in addition: an AMS \texttt{amsart} preamble that restates the article's own declarations and macros used by the retained lines, namely the environments \texttt{theorem} on separate counters, together with the standard packages those lines need.  The retained environments print in this document as Theorem 1 in order of appearance, whereas the article prints the same items with the numbering of its own full document.  The complete native source of the article as distributed in the arXiv e-print for this exact version is bundled unchanged as \texttt{sources/197657/source0.tex}.
\begin{align}
  g_{aa}(0,0)  &= 3\,\mathbb E[Df\cdot D^2f] + \mathbb E[r^*\cdot D^3f], \label{eq:gaa-exact}\\
  g_{aaa}(0,0) &= 3\,\mathbb E[(D^2f)^2] + 4\,\mathbb E[Df\cdot D^3f] + \mathbb E[r^*\cdot D^4f]. \label{eq:gaaa-exact}
\end{align}

\begin{theorem}[Near-pitchfork structure]\label{thm:near-pitchfork-structure}
Let $\sigma$ satisfy $\sigma''(0)=0$ (as holds for $\tanh$, $\mathrm{erf}$, and all odd activations).  Then:
\begin{enumerate}[label=(\roman*)]
\item $g_{aa}$ is suppressed: both terms in \eqref{eq:gaa-exact} contain the factor $D^2f$ or $D^3f$, each of which involves $h''(w_j^{*\top}x,\lambda^*)=\lambda^*\sigma''(w_j^{*\top}x)$.  Since $\sigma''(0)=0$ and $\sigma''(z)=\sigma'''(0)z+O(z^2)$, each such factor is $O(\max_j|w_j^{*\top}x|)$. Thus
\begin{equation}\label{eq:gaa-suppression}
  |g_{aa}|\le C_1\cdot\lambda^*|\sigma'''(0)|\cdot\max_j\!\sqrt{\mathbb E[(w_j^{*\top}x)^2]}\cdot\sqrt{\mathbb E[(Df)^2]\,\mathbb E[(v_0^{\top}x_{\max})^4]},
\end{equation}
where $C_1$ depends on $\|v\|$ and the data distribution but not on $m$ or $\lambda^*$ directly, and $x_{\max}$ denotes the block achieving the maximum.

\item $g_{aaa}$ is not suppressed: the dominant term $3\mathbb E[(D^2f)^2]$ involves $(h'')^2\propto (\sigma'''(0))^2(w_j^{*\top}x)^2$, which is $O(\|W^*\|^2)$ but is squared and summed, giving a quantity bounded below by
\begin{equation}\label{eq:gaaa-lower}
  g_{aaa}\ge 3(\lambda^*)^2(\sigma'''(0))^2\,\sum_j v_j^2\,\mathbb E\!\left[(w_j^{*\top}x)^2(v_{0,j}^\top x)^4\right]-C_2\cdot|\sigma'''(0)|\|W^*\|^3.
\end{equation}
\end{enumerate}

The ratio satisfies
\begin{equation}\label{eq:ratio-explicit}
  \frac{|g_{aa}|}{|g_{aaa}|} = O\!\left(\frac{1}{\lambda^*\,|\sigma'''(0)|\,\sqrt{\mathbb E[(w_{\max}^{*\top}x)^2]}}\right)
\end{equation}
as $\max_j\|w_j^*\|\to 0$. For $\sigma=\tanh$, $|\sigma'''(0)|=2$, and the ratio is small whenever the pre-activations are concentrated near zero---precisely the regime where $\lambda^*<1$.
\end{theorem}

\begin{proof}
(i) Here
\[
  D^2f=\sum_j v_j h''(w_j^{*\top}x,\lambda^*)(v_{0,j}^\top x)^2.
\]
Each summand satisfies $|h''(z,\lambda^*)|\le\lambda^*|\sigma'''(0)||z|+O(z^2)$ since $\sigma''(0)=0$.  The Gauss--Newton term $3\mathbb E[Df\cdot D^2f]$ is therefore bounded by Cauchy--Schwarz as stated.  The residual term $\mathbb E[r^*\cdot D^3f]$ contains $h'''(z,\lambda^*)=\lambda^*\sigma'''(z)$, where $\sigma'''(0)\neq 0$, so its leading contribution is $O(\lambda^*|\sigma'''(0)|\cdot\mathbb E[|r^*|\cdot|v_{0,j}^\top x|^3])$, which is small when the residual $r^*$ is small.

(ii) For $D^2f$, the square expands as
\[
  (D^2f)^2=\Bigl(\sum_j v_j h''(w_j^{*\top}x,\lambda^*)(v_{0,j}^\top x)^2\Bigr)^2,
\]
with leading term
\[
  (\lambda^*)^2(\sigma'''(0))^2\Bigl(\sum_j v_j(w_j^{*\top}x)(v_{0,j}^\top x)^2\Bigr)^2.
\]
Taking expectations and applying Jensen's inequality yields the lower bound.  The cross terms $4\mathbb E[Df\cdot D^3f]$ involve $h'\cdot h'''$; since $h'(0,\lambda^*)=1$ and $h'''(0,\lambda^*)=\lambda^*\sigma'''(0)\neq 0$, this term is
\[
  O\bigl(\lambda^*|\sigma'''(0)|\cdot\mathbb E[|Df|\cdot|v_{0,j}^\top x|^3]\bigr),
\]
which is generically smaller than the $(D^2f)^2$ term.

The ratio bound \eqref{eq:ratio-explicit} follows by dividing \eqref{eq:gaa-suppression} by \eqref{eq:gaaa-lower}.
\end{proof}

\end{document}
